
NLP benchmarks are routinely ''solved'' --- performance saturates near human parity --- yet no automated system detects when this saturation occurs, leaving benchmark retirement decisions to informal community consensus that can lag the actual event by months. This paper proposes an automated saturation detection pipeline that reframes the problem as logistic curve fitting: given a benchmark's score-over-time leaderboard timeseries, a bounded logistic growth model is fitted and formally compared against linear and power-law alternatives via the Akaike Information Criterion (AIC). Applied to GLUE and SuperGLUE, the logistic model is overwhelmingly preferred ($\Delta AIC \approx -194$ to $-250$ versus linear and $-113$ to $-357$ versus power-law alternatives --- decisive by any standard threshold), and a dual-criterion saturation detector recovers the community-recognized saturation events within 3 and 5 months respectively. A notable finding emerges: the logistic inflection point predates benchmark launch for both benchmarks ($t_0 = -6.77$ months for GLUE, $-2.85$ months for SuperGLUE), suggesting that BERT-era pretraining had already saturated benchmark-exploitable signal before public leaderboard tracking began. The pipeline also generalizes to benchmarks with as few as 30 leaderboard entries when bounded fitting is applied (mean $R^2$ = 0.913 on synthetic benchmarks). Taken together, these results establish, to the authors' knowledge, the first quantitative, reproducible foundation for automated benchmark saturation monitoring. One honest negative is reported: prospective forecasting from 6-month-early truncated data fails, as the truncated timeseries (approximately 14 months) is insufficient for logistic plateau detection, scoping the method to retrospective saturation analysis.

